\documentclass{article}

\usepackage{fancyhdr}
\usepackage[T1]{fontenc}
\usepackage[margin=0.5in]{geometry}
\usepackage{parskip}
\usepackage{amsmath}
\usepackage{amssymb}
\usepackage{parskip}

\pagestyle{fancy}
\fancyhead[l]{Andrew O'Dwyer}
\fancyhead[c]{Ritangle Question 8}
\fancyhead[r]{\today}

\begin{document}

\section{Problem Description}
\huge
A cube and a sphere have the same volume and same centre,
If p\% of the surface area of the cube lies outside the sphere,
what is p?

\section{Problem Solution}
\begin{align*}
    &\text{Let volume} = v\\
    &\text{Let side length} = s\\
    &\text{Let radius of sphere} = r_{sphere}\\
    &\text{Let radius of circle} = r_{circle}
\end{align*}

Then 

\begin{align*}
    s &= \sqrt[3]{v}\\
    r_{sphere} &= \sqrt[3]{\frac{3v}{4\pi}}\\
               &= v^\frac{1}{3} \left(\frac{3}{4\pi}\right)^\frac{1}{3}
\end{align*}

So find the area of the circle intersecting the cube face.
The radius can be found by saying the adjacent equals
$\frac{s}{2}$ and the hypotenuse equals r

\begin{align*}
    r_{circle}^2 &= r_{sphere}^2 - \frac{s}{2}^2\\
    &= v^\frac{2}{3} \left(\frac{3}{4\pi}\right)^\frac{2}{3}
    - \frac{v^\frac{2}{3}}{4}\\
    &= v^\frac{2}{3}\left(
        \left(\frac{3}{4\pi}\right)^{2/3}
        - \frac{1}{4}\right)\\
\end{align*}

So surface area of one face equals $s^2 - r_{circle}^2\pi$

\begin{align*}
    \text{So p} &= \frac{s^2 - r_{circle}^2\pi}{s^2}\\
            &= 1 - \frac{r_{circle}^2\pi}{s^2}\\
            &= 1 - \frac{v^\frac{2}{3}\left(
                \left(\frac{3}{4\pi}\right)^{2/3}
                - \frac{1}{4}\right)\pi}
                {v^{\frac{2}{3}}}\\
            &= 1 - \left(\left(\frac{3}{4\pi}\right)^{2/3}
                    - \frac{1}{4}\right)\pi\\
            &= 0.576
\end{align*}

This is all under the assumtion that the circle can fully fit
in the cubes face. Is this assumtion true?

Well for it to be true,
\begin{align*}
    &r_{circle} <= \frac{s}{2}\\
    \implies& 2\sqrt{
                    v^\frac{2}{3}\left(
                    \left(\frac{3}{4\pi}\right)^{2/3}
                    - \frac{1}{4}\right)
                } <= v^{\frac{1}{3}}\\
    \implies& 2v^\frac{1}{3}\sqrt{
                            \left(\frac{3}{4\pi}\right)^{2/3}
                            - \frac{1}{4}
                        } <= v^{\frac{1}{3}}\\
    \implies& 2\sqrt{
                \left(\frac{3}{4\pi}\right)^{2/3}
                - \frac{1}{4}
                } <= 1\\
            &\text{which is true for all v as } 0.734 < 1
\end{align*}

THE CLUE WAS 102

\end{document}          
